% Laget automatisk av elevark/lag_elevark.py fra sim/pendel-pinne.html. Endringer her blir overskrevet.
\documentclass[11pt]{article}
\usepackage{elevark}
\begin{document}
\elevarkhode{Fysikk 1 · Mekanisk energi}{Pendel med pinne}{https://fysikklab.no/sim/pendel-pinne.html}{Simuleringen}

En kule henger i en 1,0~m lang snor. Rett under opphenget står en pinne. Når pendelen svinger forbi bunnen, treffer snora pinnen, og kula svinger videre i en mindre sirkel. Hvor høyt kommer den på den andre siden?

\section*{Oppgaver}

Jobb gjerne i par. Skriv ned gjetningen og begrunnelsen \emph{før} dere trykker Start.

\begin{oppgaveliste}
\oppgave \fase{gjett}{Gjett}Pendelen slippes fra 40°, og pinnen står 0,50~m under opphenget. Kommer kula høyere, lavere eller like høyt på høyre side som der den startet?
\svar
\oppgave \fase{test}{Test}Kjør simuleringen. Flytt pinnen opp og ned og prøv igjen. Hva er likt hver gang, og hva endrer seg?
\svar
\oppgave \fase{forklar}{Forklar}Bruk energibevaring til å forklare det dere så. Hvorfor gjør ikke snorkraften noe arbeid på kula, selv om den er stor?
\svar
\oppgave \fase{test}{Regn}Regn ut farten til kula i det laveste punktet når startvinkelen er 40°. Sjekk svaret i simuleringen.
\svar
\oppgave \fase{gjett}{Gjett}Slipp pendelen fra 90° med pinnen 0,50~m under opphenget. Hva skjer? Hvorfor kommer ikke kula opp til starthøyden nå?
\svar
\oppgave \fase{utfordring}{Utfordring}Pendelen slippes fra 90°. Hvor langt ned må pinnen stå for at kula skal gå hele veien rundt pinnen med stram snor? Tips: På toppen må snorkraften være null eller større, så farten må oppfylle $v^{2} \geq gR$, der $R$ er radien rundt pinnen.
\svar
\end{oppgaveliste}

\end{document}
